\begin{center}
    \section*{\underline{Gleichstrom}}
\end{center}

\begin{tikzpicture}[remember picture, overlay]

% ---- ERSTE REIHE: Am Seitenursprung starten
\setlength{\rowheight}{22mm} % Erste Reihenhöhe: 22mm

    % A1) Erste Kachel absolut platzieren (vom Seiten-Nordwest-Anker aus)
    \setlength{\cellwidth}{30mm} % Breite einstellen
    \Tile{A1}{([xshift=10mm,yshift=-20mm]current page.north west)}{\cellwidth}{\rowheight}
        {
            \vspace{-2mm}
            \titlerm{El. Feldstärke}\\ \vspace{-3mm}
            \begin{equation*}
                \vec{F} = Q\cdot E
            \end{equation*}
            \begin{equation*}
                U = F\cdot L
            \end{equation*}
        }{0}{1}{1}{0} % Kanten: unten & rechts

    % B1) Zweite Kachel rechts anbauen
    \setlength{\cellwidth}{65mm} % Breite einstellen
    \Tile{B1}{(A1.north east)}{\cellwidth}{\rowheight}
        {
            \vspace{-2mm}
            \titlerm{Elektrische Stromstärke}\\ \vspace{1mm}
            \begin{minipage}{0.44\cellwidth}
                \centering
                \begin{equation*}
                    I = \frac{Q}{t}
                \end{equation*}
                (Bei stat. Strom)
            \end{minipage}%
            \vline
            \begin{minipage}{0.49\cellwidth}
                \centering
                \begin{equation*}
                    i = \frac{\text{d}q(t)}{\text{d}t}
                \end{equation*}
                (Bei instat. Strom)
            \end{minipage}
        }{0}{1}{1}{0} % Kanten: unten & rechts

    % C1) Dritte Kachel rechts anbauen
    \setlength{\cellwidth}{26mm} % Breite einstellen
    \Tile{C1}{(B1.north east)}{\cellwidth}{\rowheight}
        {
            \vspace{-2mm}
            \titlerm{Stromdichte}\\ \vspace{-1mm}
            \begin{equation*}
                S = \frac{I}{A}
            \end{equation*}
        }{0}{1}{1}{0} % Kante: unten & rechts
    
    % D1) Vierte Kachel rechts anbauen
    \setlength{\cellwidth}{36mm} % Breite einstellen
    \Tile{D1}{(C1.north east)}{\cellwidth}{\rowheight}
        {
            \vspace{-2mm}
            \titlerm{Ohm'sches Gesetz}\\ \vspace{-2.5mm}
            \begin{align*}
                U &= R\cdot I \\
                I &= G\cdot U
            \end{align*}
        }{0}{1}{1}{0} % Kante: unten & rechts

    % E1) Fünfte Kachel rechts anbauen
    \setlength{\cellwidth}{33mm} % Breite einstellen (Summe aller Zellen sei 190mm!)
    \Tile{E1}{(D1.north east)}{\cellwidth}{\rowheight}
        {
            \vspace{-2mm}
            \titlerm{El. Arbeit}\\ \vspace{-2.5mm}
            \begin{align*}
                W &= Q\cdot U \\
                  &= U\cdot I\cdot t
            \end{align*}
        }{0}{0}{1}{0} % Kante: unten

\end{tikzpicture}